\section{The layered variational certificate}
\label{sec:layered}

This section is a finite, computer-assisted theorem.  The analytic
equidistribution theorem proved in the preceding sections is not used in the
evaluation of the integrals.  It enters only when the certified support is
connected to the sieve in \cref{sec:main-proof}.

\subsection{The layered support}

Fix
\begin{equation}\label{eq:finite-parameters}
 k=46,\qquad
 R=\frac{521}{2000},\qquad
 \delta=\frac1{50},\qquad
 \omega=\frac{57}{10000}.
\end{equation}
The endpoint used for the marginal integral is
\begin{equation}\label{eq:outer-radius}
 A_1=\frac14+\omega=\frac{2557}{10000},\qquad
 \varepsilon=R-A_1=\frac3{625},\qquad
 R_o=A_1-\varepsilon=\frac{2509}{10000}.
\end{equation}
Let
\[
\begin{array}{c|cccccccccc}
m&1&2&3&4&5&6&7&8&9&10\\ \hline
B_m&
\frac{157}{1000}&\frac{161}{1000}&\frac{179}{1000}&
\frac{179}{1000}&\frac{188}{1000}&\frac{199}{1000}&
\frac{203}{1000}&\frac{208}{1000}&\frac{208}{1000}&
\frac{211}{1000}.
\end{array}
\]
For definiteness put $B_m=219/1000$ for $m\ge11$.

\begin{definition}[Layered support]\label{def:layered-region}
For $\mathbf t=(t_1,\ldots,t_k)\in[0,1]^k$, set
\[
 L(\mathbf t)=\{i:t_i>\delta\},
 \qquad M(\mathbf t)=|L(\mathbf t)|.
\]
The layered region is
\begin{equation}\label{eq:Omega}
 \Omega=\left\{\mathbf t\in[0,1]^k:
 \sum_{i=1}^k t_i<R,\quad
 \sum_{i\in L(\mathbf t)}t_i\le B_{M(\mathbf t)}
 \right\}.
\end{equation}
\end{definition}

Only the layers $0\le M\le10$ are nonempty.  Indeed,
$10\delta=1/5\le B_{10}$, whereas
$11\delta=11/50>B_{11}$.  The point of \cref{eq:Omega} is that it retains
both the number of coordinates above $\delta$ and their total mass.  A
radial support condition remembers neither datum.

\begin{proposition}[Exact support routing]\label{prop:support-routing}
Every pair of coordinate configurations occurring in the denominator and
the truncated marginal functionals on $\Omega$ satisfies one of the algebraic
support conditions associated with Type~I, Type~IIa, Type~IIb, Type~III,
Bombieri--Vinogradov, or the prescribed-factor condition of
\cref{thm:S52}.  In every prescribed-factor cell, the three strict
parameter inequalities in \cref{thm:S52} hold with positive rational
slack.
\end{proposition}

\begin{proof}
For fixed values of the two large-coordinate counts, all relevant
conditions are rational linear inequalities in the ordered coordinate
masses and the auxiliary partition points.  Of the
$11\cdot12/2=66$ unordered count pairs, the pair $(0,0)$ contains no large
coordinate and is routed directly to the Bombieri--Vinogradov cell.  We
enumerate admissible factor partitions for each of the remaining $65$ pairs
and ask Z3 to test the negation of their union in exact rational
quantifier-free linear arithmetic.  Every returned result is unsatisfiable.
The capacities, all covering partitions, and every strict rational margin
are recorded in the ancillary support certificate.  Thus this is a finite
exhaustive check, rather than a sampled-grid calculation.  The precise
verification interface and its trusted-kernel boundary are described in
\cref{app:certificate}.
\end{proof}

\subsection{The rational test function}

Put $x_i=t_i/R$ and define
\[
 U=\sum_i x_i,\qquad
 U_b=\sum_{i\in L(\mathbf t)}x_i,\qquad
 D=\sum_i x_i^2.
\]
For $1\le m\le10$ let
\begin{equation}\label{eq:Y-m}
 Y_m=\frac{U_b-m\delta/R}{B_m/R-m\delta/R}\in[0,1].
\end{equation}
Let $L_j(y)=P_j(2y-1)$ be the shifted Legendre polynomial.  The frozen
candidate has the form
\begin{equation}\label{eq:test-function}
 F(\mathbf t)=\1_{\Omega}(\mathbf t)
 \prod_{i=1}^{46}q(x_i)
 \left(A_{M(\mathbf t)}(U,Y_{M(\mathbf t)})
       +D\,C_{M(\mathbf t)}(U)\right),
\end{equation}
where the $Y$-dependent terms are omitted in layer $M=0$, and
\begin{equation}\label{eq:A-layer}
 A_m(U,Y)=a_{m,0}(U)+\sum_{\ell=1}^{4}a_{m,\ell}(U)L_\ell(Y).
\end{equation}
Here $q$ is the fixed positive degree-$32$ one-coordinate polynomial whose
$33$ shifted-Chebyshev coefficients, with denominator $2^{60}$, are frozen in
the ancillary profile source identified by hash in
\cref{app:certificate}.  All the $a_{m,\ell}$ and $C_m$ are expanded
in shifted Legendre polynomials in $U$.  Their degree schedules, indexed by
$m=0,\ldots,10$, are
\begin{align}
 \deg a_{m,0}&=(4,4,4,4,4,4,4,3,2,1,1),\notag\\
 \deg a_{m,\ell}\;(\ell\ge1)
 &=(-1,4,4,4,4,3,3,2,1,-1,-1),\notag\\
 \deg C_m&=(3,3,3,3,3,2,2,1,1,-1,-1),
\label{eq:degree-schedules}
\end{align}
where $-1$ means that the mode is absent.  This gives $208$ coefficients.
Every coefficient is an integer divided by $2^{50}$.  The coefficient
ordering and the complete numerator vector are part of the candidate whose
canonical SHA-256 digest is
\begin{equation}\label{eq:candidate-hash}
\begin{gathered}
\mathtt{fbc1d61a9bb8fd7576219001bee8b3c}\\
\mathtt{92728dbfa633d6e710583d8aabd8f98c9}.
\end{gathered}
\end{equation}
The digest binds the $208$ layer coefficients to $k=46$, the degree schedules,
$R_o$, and the support digest.  The polynomial $q$ is bound separately by the
profile-source hash in \cref{app:certificate}.  The floating-point searches
that found these frozen inputs are not used in their certification.

\subsection{The interval calculation}

For a square-integrable function supported on $\Omega$, define
\begin{align}
 I(F)&=\int_{[0,\infty)^k}F(\mathbf t)^2\,\dd\mathbf t,
\label{eq:I-def}\\
 J_i(F)&=
 \int_{\substack{\mathbf t_{\widehat i}\in[0,\infty)^{k-1}\\
                  \sum_{j\ne i}t_j\le R_o}}
 \left(\int_0^\infty F(\mathbf t)\,\dd t_i\right)^2
 \dd\mathbf t_{\widehat i}.
\label{eq:J-def}
\end{align}
By symmetry $J_i(F)$ is independent of $i$; denote it by $J(F)$.
After the change of variables $t_i=Rx_i$, write
\[
 I(F)=R^{46}\overline I,\qquad
 J(F)=R^{47}\overline J.
\]
Notice that \cref{eq:J-def} uses the conservative endpoint $R_o$.
No unverified contribution from its complement is added.

\begin{theorem}[Certified variational inequality]
\label{thm:finite-certificate}
The rational function \cref{eq:test-function} satisfies
\begin{align}
 \overline I&>
 9.0507319290958436555533517458740189572\cdot10^{-110},
\label{eq:I-lower}\\
 \overline J&>
 7.5531779869092232940446089973346969994\cdot10^{-111},
\label{eq:J-lower}\\
 46R\overline J-\overline I&>
 2.4125261747861770030321563214845716954\cdot10^{-114}.
\label{eq:margin-lower}
\end{align}
In particular,
\begin{equation}\label{eq:score-lower}
 \frac{46J(F)}{I(F)}>
 1.0000266555919862183480792294536290427>1.
\end{equation}
\end{theorem}

\begin{proof}
On each layer the integrands are polynomials over a rational polytope.
The exact preprocessing uses rational polynomial arithmetic; the final
integrals are evaluated at $1024$-bit precision with directed outward
rounding in Arb \cite{Arb}.  The $I$ certificate separately encloses all
eleven layer contributions.  The $J$ certificate separately encloses the
small--small, twice-small--large, and large--large channels, and it imports
the certified interval for $I$ by its file hash.  It then evaluates the
signed expression $46R\overline J-\overline I$ directly, rather than
deducing its sign by dividing two rounded decimal approximations.
The resulting interval has the positive lower endpoint displayed in
\cref{eq:margin-lower}.  Candidate and support hashes agree in all three
artifacts.  Full midpoint-radius enclosures and reproduction commands are
given in \cref{app:certificate}.
\end{proof}

\subsection{Removing the polyhedral discontinuity}

The indicator in \cref{eq:test-function} is convenient for exact
integration, whereas the sieve is usually stated for smooth test functions.
The strict margin makes the passage harmless.

\begin{lemma}[Smooth approximation]\label{lem:smooth-approximation}
For every $\epsilon>0$ there is a smooth symmetric function
$F_\epsilon$ compactly supported in the interior of $\Omega$ such that
\[
 |I(F_\epsilon)-I(F)|<\epsilon,\qquad
 |J_i(F_\epsilon)-J_i(F)|<\epsilon
 \quad(1\le i\le46).
\]
Consequently $F_\epsilon$ may be chosen so that
$46J(F_\epsilon)>I(F_\epsilon)$.
\end{lemma}

\begin{proof}
The region $\Omega$ is a finite union of bounded rational polyhedral cells
and its boundary has Lebesgue measure zero.  Thus smooth functions compactly
supported in the cell interiors are dense in $L^2(\Omega)$.  Averaging over
the symmetric group preserves the support and gives symmetric
approximants.  The map
\[
 T_iG(\mathbf t_{\widehat i})=\int_0^\infty G(\mathbf t)\,\dd t_i
\]
has operator norm at most $R^{1/2}$ from $L^2(\Omega)$ to
$L^2(\dd\mathbf t_{\widehat i})$, by Cauchy--Schwarz on each fibre.
Restriction to $\sum_{j\ne i}t_j\le R_o$ cannot increase the norm.
Hence $I(G)=\|G\|_2^2$ and
$J_i(G)=\|\1_{\{\sum_{j\ne i}t_j\le R_o\}}T_iG\|_2^2$
are continuous under $L^2$ approximation.  The positive margin in
\cref{eq:margin-lower} therefore persists.
\end{proof}

Together, \cref{prop:support-routing,thm:finite-certificate,lem:smooth-approximation}
prove \cref{thm:finite-intro}.
